\chapter*{Preface}
\addcontentsline{toc}{chapter}{Preface}

This book accompanies the course \emph{Advanced Algorithms} at Utrecht
University, and it has no other purpose. It is not a reference work, not a
research monograph, and not an attempt to survey a field. It is the book for
one course, written for the people taking it, and that single purpose is the
reason it can be as short as it is: a topic is in this book if the course
covers it, and out if it does not, however interesting.

\section*{Who it is for}

It is written for students who have had a first algorithms course. We assume
asymptotic notation, the standard graph algorithms, dynamic programming, and
the notions of \textbf{P}, \textbf{NP} and \textrm{NP}-completeness; everything
beyond that is developed here. \Cref{app:preliminaries} says precisely what is
assumed, and it opens with ten questions that take about a quarter of an hour
and tell you which of its sections you need to read --- including the case
where the honest answer is that a prerequisite is missing, which it treats as a
question of sequencing rather than of ability, and answers with routes rather
than a verdict.

The course draws people from several degree programmes and the book is written
in the knowledge that the mathematical background in the room is uneven.
Nothing here needs mathematics beyond a bachelor's degree in a quantitative
subject. What it does need, and cannot supply, is the willingness to read a
proof slowly.

\section*{How this book was written}

This book was drafted by Claude, an AI system made by Anthropic, working from
the lecture slides and notes of the course under the direction of its authors.
The instructions it worked to are in the repository, and the short version is:
every chapter was written against a specific source deck and checked against it
in both directions, so that nothing lectured is missing and nothing is here that
the lectures do not cover; every running time quoted from a slide was verified
against the literature rather than copied, because several had gone stale and
one was wrong; and every \textsc{doi} in the bibliography was checked against
the publisher's record rather than written from memory.

That process catches a great deal and it does not catch everything. A
mathematical review of the second half of the book found a clause gadget that
did not work, a reduction quoted in the wrong direction, an algorithm
attributed to the wrong paper, a worked example whose arithmetic was impossible,
and a factor of two --- all of which are now fixed, and all of which were the
kind of error that survives a confident-sounding paragraph. The reader should
expect that others remain, and should read accordingly: the proofs in this book
are meant to be checked, not believed. Where something looks wrong, it may
be.

We think the trade is worth making. A course book that exists and is being
corrected is more use than a better one that never gets written, and this one
was assembled in a fraction of the time it would otherwise have taken. But the
reader is entitled to know how it was made, which is why this section is here
rather than in a footnote.

\todo{Add the errata channel: where should a student send a mistake? An email
address, a GitHub issue on the course repository, or a note to the teaching
assistants --- whichever you prefer, but it should be named here in one
sentence, because the section above asks people to report errors and then does
not say how.}

\section*{How to use it}

The course is examined with the
\href{https://openpoolexams.science.uu.nl/}{open pool exam method}: the pool of
possible exam questions is published in advance, and at least $30\%$ of the
questions on the midterm and the final exam are taken \emph{verbatim} from that
pool. The pool is reproduced in \cref{app:pool}.

This shapes the book more than anything else about it. Rather than keep two
lists that would drift apart, the book keeps one: the learning goals at the head
of each chapter \emph{are} the pool questions, and a goal marked \pooltag{} can
be set on the exam word for word. Everything needed to answer it is in the body
of that chapter. Where a section goes beyond what the lectures cover, its
heading says \emph{(non-examinable)} and a note underneath says why it earns its
place; nothing marked that way is needed for the pool.

If you can answer every pool question using only this book, you have understood
the course. The pool is a floor and not a ceiling: the remaining exam questions
test the same material on problems you have not seen before, which is what the
exercises at the end of each chapter are for.

\Cref{ch:introduction} says what the subject is and how the five parts fit
together, and is the place to start.

\section*{Sources and acknowledgements}

This book did not start from nothing. The course \emph{Advanced Algorithms} was
built by \textbf{Jesper Nederlof} and \textbf{Johan van Rooij}, and most of it
is based on their lecture slides: the selection of topics, the order they come
in, and a good number of the individual arguments are theirs. Each chapter's
``Notes and further reading'' says which material it follows.

Part~II --- real models of computation, the existential theory of the reals,
and its geometric consequences --- is based on lecture notes by
\textbf{Tillmann Miltzow}.

\todo{Still to acknowledge: the teaching assistants, and the students who
report errata --- the second list is worth keeping from the first year onwards,
because naming people who found mistakes is the cheapest way to get more of
them found.}
